%! Sinusoidal Function Transformations — Unit 6, Lesson 6.7 — Exit Ticket
\documentclass[10pt]{article}
\usepackage{precalculus-article}
\usepackage{precalculus-boxes}
\newcommand{\TallMath}[1]{$\displaystyle #1\rule[-1.4em]{0pt}{3.2em}$}

\begin{document}

\pageheader{Unit 6, Lesson 6.7}{Exit Ticket}
\namedateperiod

\vspace{0.12in}

\begin{notesbox}{Exit Ticket --- Answer independently}

\textbf{1.}\quad Identify the amplitude, period, phase shift, and midline of the function
$$f(\theta) = -3\sin\!\left(2\!\left(\theta + \dfrac{\pi}{4}\right)\right) - 1$$

\vspace{4pt}
\renewcommand{\arraystretch}{1.6}
\begin{tabular}{ll}
Amplitude: & \underline{\hspace{4cm}} \\[4pt]
Period: & \underline{\hspace{4cm}} \\[4pt]
Phase shift: & \underline{\hspace{4cm}} \\[4pt]
Midline: & \underline{\hspace{4cm}} \\
\end{tabular}

\vspace{0.16in}
\textbf{2.}\quad Circle the letter of the equation that matches this description:

\textit{``A sinusoidal function with amplitude 4, period $\pi$, no phase shift, and midline $y = 2$.''}

\vspace{6pt}
\begin{tabular}{ll}
(A) $y = 4\sin(\pi\theta) + 2$ \qquad\qquad & (B) $y = 4\sin(2\theta) + 2$ \\[6pt]
(C) $y = 2\sin(4\theta) + \pi$ & (D) $y = 4\sin(\theta) + 2$ \\
\end{tabular}

\vspace{0.16in}
\textbf{3.}\quad In one sentence, explain the difference between the \emph{period}
and the \emph{amplitude} of a sinusoidal function.

\writelines{2}

\end{notesbox}

\end{document}
